\section{Related Work}
\label{sec:related}

\textbf{Evaluation harnesses.} Large language model evaluation is increasingly harness-driven: HELM standardizes scenarios and metrics~\citep{liang2022helm}, \texttt{lm-evaluation-harness} provides a reproducible task runner~\citep{gao2023lm}, and BIG-bench coordinates many contributors around a fixed interface~\citep{srivastava2022bigbench}. These harnesses standardize \emph{task} execution. Our problem is different: a medical VLM study has a small number of tasks but a long, heterogeneous execution path (image preprocessing, model-specific processors, structured parsing, contract checks), and the dominant failure modes are not task coverage but silent data corruption and incomplete resumption.

\textbf{Reproducibility instruments.} Datasheets~\citep{gebru2021datasheets} and model cards~\citep{mitchell2019model} document data and models; foundation-model reporting has made provenance a first-class concern~\citep{bommasani2021foundation}. Our manifest lock, coverage census, and per-stage state file are task-specific analogues: they document the \emph{evaluation}, not just the dataset or the model.

\textbf{Medical VLM evaluation and benchmarks.} Medical VLM evaluation has moved from single-task accuracy to robustness under artefacts~\citep{cheng2025robustness}, prompt and workflow stress~\citep{yang2026stress}, visual grounding~\citep{zafar2026beyond}, calibration/abstention~\citep{qu2026accuracy}, triage~\citep{li2025atcxr}, and transfer failure~\citep{fisher2025negative}. These studies rely on the same execution primitives---batching, parsing, split filtering---but rarely report them. Benchmarks such as OmniMedVQA~\citep{hu2023omnimedvqa} and chest X-ray corpora~\citep{wang2017chestxray8,bustos2020padchest,irvin2019chexpert,johnson2019mimic} define tasks and labels, but their existence does not solve the execution problem: the same dataset can yield different denominators under different split filters, label encodings, and parser conventions, and public datasets carry known label-quality and domain-shift limits~\citep{rafferty2026limitations}. The study uses CheXagent~\citep{chen2024chexagent} and MedGemma~\citep{sellergren2025medgemma} as two of its models; their model-specific processors are part of the measurement declaration. We do not re-evaluate the scientific claims of these works; we evaluate whether the harness that produced them is auditable.

\textbf{Where our baselines come from.} Four families bear on this problem and none of them is compared against in the medical VLM literature. Multimodal evaluation kits such as VLMEvalKit reuse a prediction file when it exists~\citep{duan2024vlmevalkit}, which is file-level resumability: correct if the file is complete, silent if it is not. Workflow engines --- Make~\citep{feldman1979make}, Snakemake~\citep{molder2021snakemake} and DVC~\citep{dvc2024} --- resume by content hash, which is strictly stronger than a stage state file, but they treat a model call as an opaque job and know nothing about output contracts. Agent-oriented evaluation frameworks such as Inspect AI log every sample and enforce retry and error policies~\citep{inspect2024}, which is closest to the row-level guarantee here. Databases solved the same durability problem decades ago, and SQLite's write-ahead log is a fully general answer available in the standard library~\citep{sqlitewal2024}. Fault injection itself is an established technique~\citep{hsueh1997fault}, and dataset metadata formats such as Croissant~\citep{akhtar2024croissant} address the adjacency of a machine-readable artifact contract. Our comparison uses one concrete instance of each family --- naive append, append plus fsync, atomic whole-file rewrite, and SQLite --- so that the durability claim has a reference class rather than a self-referential test.

\textbf{Durable execution.} Write-ahead logging and atomic replace are standard database techniques~\citep{mohan1992aries}; we apply the same ideas to append-only JSONL prediction files. We do not claim the mechanism as novel. The design target is not transactional throughput but a single long-running GPU job on a desktop machine, where a power cut or driver reset must cost at most the in-flight batch --- and the measured question is whether the mechanism is worth its 1.95\,ms per cell.

\textbf{Technical debt in ML systems.} Sculley et al.~\citep{sculley2015hidden} describe how data dependencies, pipeline jungles, and configuration debt accumulate in production ML; Breck et al.~\citep{breck2017mltest} turn production readiness into a test rubric. Our harness is a benchmark-side instantiation of the same concern: the evaluation pipeline is a system with its own failure modes, and those failure modes need tests.
